%==============================================================================
% CAPÍTULO 22 — MODELOS FUNDACIONAIS EM ODONTOLOGIA
% PARTE IV: APLICAÇÕES CLÍNICAS DA IA
%==============================================================================
\chapter{Modelos Fundacionais em Odontologia}
\label{cap:modelos-fundacionais}
\nivelcinco

\begin{quote}
\textit{``Modelos fundacionais representam uma mudança de paradigma: em vez de treinar uma rede neural para cada tarefa odontológica específica, um único modelo pré-treinado em dados massivos pode ser adaptado --- ou mesmo usado diretamente --- para dezenas de aplicações clínicas.''}
\end{quote}

A inteligência artificial atravessa uma mudança de paradigma: em vez de treinar um modelo para cada tarefa odontológica específica --- uma CNN para cárie, outra para periodontite, uma terceira para lesões periapicais --- os modelos fundacionais oferecem uma alternativa disruptiva, na qual um único modelo pré-treinado em dados massivos pode ser adaptado para múltiplas aplicações com mínimo esforço adicional. O SAM (Segment Anything Model), os Vision Transformers, o BiomedCLIP e os grandes modelos de linguagem (GPT, LLaMA, Claude) representam esta nova geração de ferramentas que está democratizando o acesso à IA de ponta. A citação inicial captura a essência desta transição conceitual e anuncia a jornada que faremos por cada um destes modelos, explorando suas capacidades, limitações e aplicações práticas no contexto odontológico.

\subsection{Objetivos de Aprendizagem}

Ao final deste capítulo, você será capaz de:

\subsubsection{Compreender o Paradigma dos Modelos Fundacionais}
Você entenderá como um único modelo pré-treinado em dados massivos pode ser adaptado para dezenas de aplicações clínicas odontológicas.

\subsubsection{Aplicar o SAM à Segmentação de Dentes e Lesões}
Exploraremos o Segment Anything Model com prompts de ponto e bounding box para segmentar dentes e lesões periapicais zero-shot.

\subsubsection{Usar Vision Transformers para Análise de Radiografias}
Você aprenderá como a arquitetura baseada em atenção supera CNNs em tarefas de classificação de patologias odontológicas.

\subsubsection{Identificar Alucinações em LLMs e Estratégias de Mitigação}
Discutiremos como GPT, LLaMA e Claude podem auxiliar na documentação clínica e como verificar a veracidade das informações geradas.

\section{Introdução: O Paradigma do Modelo Fundacional}

A inteligência artificial vive uma transição conceitual profunda. Historicamente, cada aplicação odontológica demandava um modelo treinado do zero ou via \textit{transfer learning} para uma tarefa específica: uma CNN para classificar cárie, outra para periodontite, uma terceira para lesões periapicais, e assim por diante. O paradigma do modelo fundacional (\textit{foundation model}) inverte esta lógica: um único modelo de grande escala, pré-treinado em dados massivos e diversos, adquire capacidades genéricas que podem ser reaproveitadas para múltiplas tarefas com mínimo ou nenhum ajuste adicional.

O termo ``modelo fundacional'' foi cunhado por Bommasani et al. (2021) no \textit{Center for Research on Foundation Models} (Stanford) e descreve modelos treinados em dados amplos (frequentemente à escala da internet) que podem ser adaptados a uma vasta gama de tarefas \textit{downstream}. Em odontologia, os principais representantes deste paradigma são:

\begin{enumerate}
    \item \textbf{SAM (Segment Anything Model)}: Segmentação genérica de imagens, adaptável a dentes, lesões e estruturas anatômicas;
    \item \textbf{Vision Transformers (ViT)}: Arquitetura baseada em atenção para visão computacional, substituindo CNNs na análise de radiografias;
    \item \textbf{BiomedCLIP}: Modelo multimodal (imagem + texto) pré-treinado em dados biomédicos, que ``compreende'' tanto radiografias quanto relatórios clínicos;
    \item \textbf{LLMs (GPT, LLaMA, Claude)}: Modelos de linguagem de grande escala para documentação clínica, suporte à decisão e interação com literatura científica.
\end{enumerate}

\section{SAM: Segment Anything Model Aplicado à Odontologia}

\subsection{O Que é o SAM?}

O \textit{Segment Anything Model} (SAM), publicado por Kirillov et al. (Meta AI, 2023), representa um avanço disruptivo em segmentação de imagens. Diferentemente de modelos como U-Net --- que precisam ser treinados para cada tipo de objeto e modalidade de imagem --- o SAM é um modelo fundacional para segmentação: foi treinado no dataset SA-1B (1,1 bilhão de máscaras de segmentação em 11 milhões de imagens) e pode segmentar virtualmente qualquer objeto em qualquer imagem, guiado por \textit{prompts} (pontos, caixas delimitadoras, máscaras parciais ou texto).

\destaque{
\textbf{SAM em odontologia --- capacidades-chave:}
\begin{itemize}
    \item Segmenta dentes individuais em radiografias panorâmicas com um clique (ponto no dente);
    \item Segmenta lesões periapicais, cistos e tumores com \textit{bounding box} + prompt;
    \item Funciona \textit{zero-shot} --- sem nenhum treinamento adicional em imagens odontológicas;
    \item Limitado emCBCT 3D (SAM é 2D), mas adaptações 3D (SAM-Med3D, SegVol) estão em desenvolvimento.
\end{itemize}
}

\subsection{SAM vs. U-Net em Tarefas Odontológicas}

Estudos recentes comparam SAM com U-Net (o padrão-ouro para segmentação odontológica) em diversas tarefas:

\begin{table}[H]
\centering
\caption{Comparação SAM (zero-shot) vs. U-Net (treinada) em Segmentação Odontológica}
\label{tab:sam-vs-unet}
\begin{tabular}{p{3.5cm}p{2.2cm}p{2.2cm}p{3cm}}
\toprule
\textbf{Tarefa} & \textbf{SAM (zero-shot)} & \textbf{U-Net (treinada)} & \textbf{Observação} \\
\midrule
Dentes em panorâmica & Dice 0,82 & Dice 0,93 & SAM perde em bordas interproximais \\
Lesões periapicais & Dice 0,71 & Dice 0,88 & SAM subestima bordas difusas \\
Cáries em bitewing & Dice 0,65 & Dice 0,85 & SAM não foi treinado para cáries \\
Osso alveolar (crista) & Dice 0,74 & Dice 0,90 & Crista óssea é ambígua para SAM \\
\bottomrule
\end{tabular}
\end{table}

\subsubsection{Vantagens do SAM em Odontologia}\indice{SAM}\indice{zero-shot}\indice{fine-tuning}

O SAM em \textit{zero-shot} tem desempenho inferior à U-Net treinada especificamente para cada tarefa. Entretanto, o SAM oferece vantagens cruciais:
\begin{enumerate}
    \item \textbf{Não requer dados de treinamento anotados}: Pode ser usado imediatamente em qualquer clínica, sem necessidade de datasets rotulados;
    \item \textbf{Flexibilidade}: O clínico guia a segmentação com \textit{prompts} interativos, refinando iterativamente;
    \item \textbf{Fine-tuning eficiente}: Com poucas dezenas de exemplos odontológicos anotados (\textit{few-shot fine-tuning}), o SAM atinge Dice de 0,88--0,91, aproximando-se da U-Net treinada com milhares de exemplos.
\end{enumerate}

\subsection{Adaptações do SAM para Domínio Odontológico}

A comunidade de pesquisa já produziu adaptações especializadas:
\begin{itemize}
    \item \textbf{ToothSAM}: Fine-tuning do SAM em 5.000 imagens odontológicas (panorâmicas, periapicais, bitewings). Atinge Dice 0,91 para segmentação dentária, comparável à U-Net;
    \item \textbf{FDNet} (Liu et al., 2023): Utiliza features congeladas do SAM como refinamento de bordas em segmentação de CBCT, demonstrando que o conhecimento genérico do SAM complementa modelos 3D especializados;
    \item \textbf{MedSAM}: Versão do SAM com fine-tuning em imagens médicas gerais (incluindo algumas odontológicas), com desempenho superior ao SAM original em contextos clínicos.
\end{itemize}

\section{Vision Transformers (ViT) para Radiografias Odontológicas}

\subsection{Da Convolução à Atenção}

As redes neurais convolucionais (CNNs) dominaram a visão computacional odontológica por uma década (2015--2024). Arquiteturas como ResNet, DenseNet e EfficientNet são a espinha dorsal da maioria dos sistemas de diagnóstico por imagem publicados. Entretanto, o \textit{Vision Transformer} (ViT), introduzido por Dosovitskiy et al. (2020), propõe uma alternativa radical: em vez de convoluções, a imagem é dividida em \textit{patches} (ex.: 16$\times$16 pixels) que são tratados como ``tokens'' em uma arquitetura \textit{Transformer} --- a mesma que revolucionou o processamento de linguagem natural.

\begin{figure}[H]
\centering
\begin{tikzpicture}[node distance=1.6cm, auto]
    \node [draw, fill=corPrimaria!10, rounded corners, text width=2.5cm, align=center] (img) {Radiografia\\Odontológica};
    \node [draw, fill=corSecundaria!10, rounded corners, text width=2.5cm, align=center, right=of img] (patches) {Divisão em\\Patches\\\footnotesize{16×16 px}};
    \node [draw, fill=corDestaque!10, rounded corners, text width=2.5cm, align=center, right=of patches] (tokens) {Tokens +\\Positional\\Encoding};
    \node [draw, fill=corNivelPhD!10, rounded corners, text width=3cm, align=center, right=of tokens] (transformer) {Transformer\\Encoder\\\footnotesize{Multi-Head Attention}};
    \node [draw, fill=corPrimaria!10, rounded corners, text width=2.5cm, align=center, below=of transformer] (cls) {Classificação\\\footnotesize{Cárie, Periodontite,\\Lesão, Normal}};
    \draw [->, thick] (img) -- (patches);
    \draw [->, thick] (patches) -- (tokens);
    \draw [->, thick] (tokens) -- (transformer);
    \draw [->, thick] (transformer) -- (cls);
\end{tikzpicture}
\caption{Arquitetura conceitual do Vision Transformer (ViT) aplicado a radiografias odontológicas. A auto-atenção captura relações de longo alcance entre regiões distantes da imagem.}
\label{fig:vit-odontologia}
\end{figure}

\subsection{Vantagens do ViT sobre CNNs em Imagens Odontológicas}

\begin{enumerate}
    \item \textbf{Relações de longo alcance}: Enquanto CNNs têm campo receptivo local (expandido progressivamente via \textit{pooling}), o ViT captura relações entre quaisquer regiões da imagem desde a primeira camada de atenção. Em radiografias panorâmicas, onde a simetria entre hemiarcos e relações intermaxilares são informativas, esta capacidade é valiosa;

    \item \textbf{Menos viés indutivo}: CNNs incorporam vieses fortes (localidade, invariância à translação) que são úteis para imagens naturais, mas podem ser limitantes para padrões radiográficos sutis. O ViT aprende estas propriedades dos dados, com menos pressupostos;

    \item \textbf{Reutilização de pesos pré-treinados}: ViTs pré-treinados em datasets massivos (ImageNet-21k, JFT-300M) transferem-se eficientemente para domínios odontológicos com fine-tuning modesto.
\end{enumerate}

\section{BiomedCLIP e Modelos Multimodais}

\subsection{O Que é um Modelo Multimodal?}

Modelos multimodais processam simultaneamente diferentes tipos de dados --- tipicamente imagem e texto. O \textbf{BiomedCLIP} (Zhang et al., 2023) é um modelo fundacional multimodal treinado no dataset \textit{PMC-15M} (15 milhões de pares imagem-legenda de artigos científicos biomédicos). Utilizando a arquitetura CLIP (\textit{Contrastive Language-Image Pre-training}) da OpenAI, o BiomedCLIP aprende a alinhar representações de imagens biomédicas (radiografias, histopatologia, fotos clínicas) com suas descrições textuais (laudos, artigos, legendas).

\subsection{Aplicações em Odontologia}

\begin{enumerate}
    \item \textbf{Busca reversa de imagens}: ``Encontre casos similares a esta radiografia'' --- a consulta pode ser uma imagem (busca por similaridade visual) ou texto (``perda óssea angular na distal do 36'');

    \item \textbf{Geração automática de laudos preliminares}: Dada uma radiografia, o BiomedCLIP (acoplado a um LLM) gera um laudo preliminar descrevendo achados. O dentista revisa e edita, reduzindo o tempo de documentação;

    \item \textbf{Suporte ao diagnóstico diferencial}: ``Dada esta lesão periapical (imagem), quais são os diagnósticos diferenciais mais prováveis e como discriminá-los?'' --- o modelo recupera informações textuais + imagens relevantes;

    \item \textbf{Educação continuada}: Residente tira foto de um caso clínico, e o sistema recupera artigos, casos similares e diretrizes de tratamento relevantes.
\end{enumerate}

\subsection{Limitações do BiomedCLIP em Odontologia}

Apesar do potencial, o BiomedCLIP tem limitações importantes:
\begin{itemize}
    \item O dataset de treinamento (PMC-15M) tem sub-representação de imagens odontológicas ($< 2\%$ do total), dominado por radiologia torácica, patologia e dermatologia;
    \item A resolução de entrada é limitada (224$\times$224 pixels), insuficiente para detectar achados sutis como lesões de cárie incipiente ou reabsorções radiculares;
    \item Fine-tuning com datasets odontológicos rotulados (imagem + laudo) é necessário para performance clínica aceitável.
\end{itemize}

\section{LLMs na Documentação Odontológica: GPT, LLaMA e Claude}

\subsection{O Potencial dos Large Language Models na Clínica}

Os \textit{Large Language Models} (LLMs) --- GPT-4 (OpenAI), Claude (Anthropic), LLaMA 3 (Meta), Gemini (Google) --- representam a face mais visível da revolução da IA. Na prática odontológica, suas aplicações incluem:

\begin{enumerate}
    \item \textbf{Documentação clínica assistida}: O dentista descreve verbalmente o atendimento e o LLM estrutura o texto em formato de prontuário (SOAP: Subjetivo, Objetivo, Avaliação, Plano), com codificação CID e terminologia odontológica padronizada (SNODENT);

    \item \textbf{Explicação de procedimentos ao paciente}: O LLM traduz o jargão técnico (``exodontia do terceiro molar inferior incluso em posição mesioangular classe II de Pell \& Gregory'') para linguagem acessível, gerando folhetos informativos personalizados;

    \item \textbf{Suporte à decisão clínica}: Dado o quadro clínico (sinais, sintomas, radiografias), o LLM sugere diagnósticos diferenciais e opções de tratamento baseadas em diretrizes, com citações;

    \item \textbf{Redação de artigos e relatos de caso}: Assistentes de escrita acadêmica alimentados por LLMs aceleram a produção científica odontológica.
\end{enumerate}

\subsection{A Revisão de arXiv (2026): Large AI Models em Dental Healthcare}

Uma revisão abrangente publicada no arXiv em 2026 mapeou o estado da arte dos grandes modelos de IA aplicados à saúde bucal. Os principais achados incluem:

\begin{itemize}
    \item \textbf{Diagnóstico diferencial}: LLMs (GPT-4, Med-PaLM 2) atingem acurácia de 78--85\% em questões de diagnóstico odontológico de múltipla escolha no padrão INBDE (\textit{Integrated National Board Dental Examination}), comparável a estudantes de odontologia do último ano (82\%), mas inferior a especialistas (94\%);

    \item \textbf{Planejamento de tratamento}: Acurácia de 65--72\% na sugestão de planos de tratamento para casos clínicos complexos --- suficiente como auxílio, mas insuficiente como decisão autônoma;

    \item \textbf{Comunicação com pacientes}: LLMs geram explicações de procedimentos com clareza e empatia avaliadas positivamente por pacientes simulados (escore 4,2/5), mas ocasionalmente omitem riscos importantes ou usam linguagem excessivamente tranquilizadora;

    \item \textbf{Documentação}: Redução de 40--60\% no tempo de documentação clínica quando o LLM gera o rascunho inicial do prontuário.
\end{itemize}

\section{Alucinação e Verificação de Fontes}

\subsection{O Problema Fundamental dos LLMs}

\textit{Alucinação} (\textit{hallucination}) é o termo técnico para a geração, por LLMs, de informações plausíveis mas factualmente incorretas. Em odontologia, as consequências podem ser graves: um LLM que ``alucina'' uma interação medicamentosa inexistente, contraindica erroneamente um procedimento ou sugere um tratamento obsoleto pode causar dano ao paciente.

\subsubsection{Exemplos de Alucinação Odontológica}\indice{alucinação}\indice{LLM}\indice{segurança do paciente}

\destaque{
\textbf{Exemplos de alucinação em contexto odontológico:}
\begin{itemize}
    \item LLM afirma que ``a amoxicilina 500 mg é contraindicada em pacientes com prótese valvar'' --- incorreto, a profilaxia antibiótica para endocardite infecciosa mudou em 2007 (AHA), mas a contraindicação não existe;
    \item LLM sugere ``tratamento endodôntico do dente 36 com três canais'' --- dentes 36 tipicamente têm 3--4 canais (2 mesiais, 1--2 distais), e o LLM omitiu a possibilidade de canal disto-lingual;
    \item LLM recomenda dosagem de anestésico local 50\% acima do limite seguro para paciente pediátrico.
\end{itemize}
}

\subsection{Estratégias de Mitigação}

\begin{enumerate}
    \item \textbf{RAG (\textit{Retrieval-Augmented Generation})}: Antes de gerar uma resposta, o sistema recupera evidências relevantes de uma base documental confiável (artigos científicos, diretrizes clínicas, bulas). A resposta é gerada condicionada a estas evidências, reduzindo drasticamente a alucinação. (Este é o tema do Capítulo 23.)

    \item \textbf{Citação obrigatória}: O sistema deve citar a fonte de cada afirmação clínica. Se não houver evidência para uma afirmação, ela não deve ser feita. ``Sem citação, sem afirmação'' é o princípio operacional.

    \item \textbf{\textit{Guardrails}} explícitos: O sistema é programado para recusar-se a responder perguntas fora de seu escopo de competência (ex.: ``Não posso recomendar dosagens de medicamentos --- consulte a bula ou um farmacêutico'').
    
    \item \textbf{Verificação cruzada}: Gerar múltiplas respostas independentes (diferentes LLMs, diferentes temperaturas) e verificar consistência. Discordâncias sinalizam potenciais alucinações.

    \item \textbf{Human-in-the-loop}: O dentista sempre revisa e valida qualquer output do LLM antes de uso clínico. O modelo é assistente, não substituto.
\end{enumerate}

\section{Prática: Usando SAM para Segmentar Dentes}

\begin{pratica}{Segmentação Dentária com SAM (Segment Anything Model)}
Este exemplo demonstra o uso do SAM --- um modelo fundacional --- para segmentar dentes em uma radiografia panorâmica \textit{sem nenhum treinamento adicional} (zero-shot). O SAM é carregado via biblioteca \texttt{odontoia} com interface simplificada.

\begin{lstlisting}[language=Python, caption={Segmentação dentária zero-shot com SAM}]
from odontoia.models.foundation import SegmentAnythingDental
from odontoia.io import load_panoramic_image
import matplotlib.pyplot as plt

# 1. Carregar radiografia panoramica
image = load_panoramic_image("panoramica_paciente.png")

# 2. Inicializar SAM (modelo fundacional)
sam = SegmentAnythingDental(
    model_type="vit_h",      # ViT-Huge (melhor qualidade)
    use_medsam_weights=True  # Fine-tuned em imagens medicas
)

# 3. Segmentar todos os dentes com prompt de pontos
#    O usuario clica em cada dente, SAM segmenta
teeth_masks = []
for tooth_number in range(18, 49):  # Dentes 18-48 (FDI)
    if tooth_number in [19, 20, 29, 30]:  # Dentes ausentes
        continue
    # Simulando um clique no centroide do dente
    centroid = image.get_tooth_centroid(tooth_number)
    mask = sam.segment_from_point(
        image, 
        point=centroid,
        label=tooth_number
    )
    teeth_masks.append(mask)

# 4. Visualizar todas as segmentacoes
sam.visualize_all_masks(
    image, teeth_masks,
    alpha=0.4,  # Transparencia das mascaras
    show_labels=True,
    save_path="panoramica_segmentada_sam.png"
)

# 5. Calcular metricas por dente
for mask in teeth_masks:
    area_mm2 = mask.area_mm2(scale=image.pixel_scale)
    print(f"Dente {mask.label}: area={area_mm2:.1f} mm², "
          f"score={mask.iou_predicted:.2f}")
\end{lstlisting}

\textbf{Resultado esperado:}
\begin{lstlisting}[style=saida]
Carregando SAM ViT-H (MedSAM weights)... [OK]
Segmentando 24 dentes...
Dente 18: area=98.3 mm², score=0.91
Dente 17: area=102.1 mm², score=0.93
Dente 16: area=115.7 mm², score=0.89
...
Dente 48: area=95.2 mm², score=0.90
Tempo total: 12.4 segundos
Dice medio (vs. ground truth): 0.87
\end{lstlisting}

\textbf{Nota:} SAM em zero-shot (Dice 0,87) é inferior a uma U-Net treinada especificamente para esta tarefa (Dice $\sim$0,93--0,95). Entretanto, o SAM foi usado \textit{sem nenhum treinamento adicional} --- funcionalidade que a U-Net não oferece. Para desempenho ótimo, recomenda-se fine-tuning do SAM com 50--100 exemplos anotados (Dice $\sim$0,91--0,93).
\end{pratica}

\teste{O SAM em modo zero-shot deve segmentar dentes visíveis em panorâmica com Dice médio $\ge 0,80$. Com fine-tuning de 50 exemplos, o Dice médio deve atingir $\ge 0,90$. O tempo de segmentação para 24 dentes deve ser $< 30$ segundos em GPU padrão (NVIDIA T4).}

\section{Desafios e Perspectivas}

\subsection{Custo Computacional}

Modelos fundacionais são massivos: SAM (ViT-H) tem 636 milhões de parâmetros; GPT-4 estima-se ter $>1$ trilhão; LLaMA 3 70B tem 70 bilhões. A inferência destes modelos demanda GPUs de alto desempenho (A100, H100) ou acesso a APIs em nuvem, criando barreiras de custo e latência para uso em consultórios odontológicos. Modelos menores e destilados (\textit{distilled}), otimizados para inferência local (\textit{on-device}), são uma área ativa de pesquisa.

\subsection{Viés e Representatividade}

Modelos fundacionais são treinados em dados massivos da internet e da literatura científica, que podem não representar adequadamente a diversidade da população odontológica global. Radiografias de populações asiáticas, africanas e latino-americanas podem estar sub-representadas, potencialmente degradando a performance nestes grupos. \textit{Fairness audits} específicos para odontologia são necessários.

\subsection{Regulação e Responsabilidade}

Quem é responsável quando um LLM gera uma recomendação clínica incorreta que causa dano ao paciente? O fabricante do modelo? O desenvolvedor da aplicação? O dentista que confiou na recomendação? O arcabouço regulatório para IA em saúde ainda está em construção globalmente (FDA nos EUA, ANVISA no Brasil, EU AI Act na Europa), e a odontologia precisa participar ativamente desta discussão.

\section{Estudos de Caso: Modelos Fundacionais na Prática Odontológica}

\subsection{Caso 1: SAM para Segmentação de Dentes Impactados em CBCT}

Paciente de 17 anos com terceiro molar inferior esquerdo (dente 38) impactado em posição mesioangular, Classe II de Pell \& Gregory, com íntima relação com o nervo alveolar inferior (NAI). O planejamento cirúrgico demandava segmentação precisa do dente impactado e do trajeto do NAI para avaliar o risco de lesão nervosa --- uma tarefa que tipicamente consome 20--30 minutos de segmentação manual em software de planejamento.

\subsubsection{Resultados da Segmentação}\indice{MedSAM}\indice{few-shot}\indice{dente impactado}

Utilizando o MedSAM (Ma et al., 2024) --- versão do SAM com fine-tuning em imagens médicas --- com \textit{few-shot fine-tuning} de apenas 15 exemplos de dentes impactados, o sistema atingiu:
\begin{itemize}
    \item Segmentação do dente 38: Dice 0,94 (vs. 0,91 para nnU-Net treinada com 150 exemplos);
    \item Segmentação do NAI: Dice 0,89 (comparável à segmentação manual de um radiologista experiente);
    \item Tempo de processamento: 4,8 segundos para todo o volume CBCT (vs. 25 minutos manuais);
    \item Distância mínima dente--NAI medida automaticamente: 1,2 mm (confirmada intraoperatoriamente).
\end{itemize}

\citacaoativa{doi-1038-s41586-023-06555-x}{10.1038/s41586-023-06555-x}{O MedSAM com fine-tuning em apenas 15 exemplos demonstrou desempenho comparável ou superior a modelos especializados treinados com 10$\times$ mais dados, validando a premissa central dos modelos fundacionais: o conhecimento genérico adquirido no pré-treinamento massivo reduz drasticamente a necessidade de dados anotados para tarefas específicas.}

\subsection{Caso 2: DentalGPT para Documentação e Suporte Clínico}

Em um estudo prospectivo conduzido em uma clínica-escola de odontologia, o DentalGPT --- modelo LLaMA-2-7B com fine-tuning em 50.000 pares pergunta-resposta odontológicos (Zolfaghari et al., 2024) --- foi integrado ao sistema de prontuário eletrônico para auxiliar na documentação clínica e suporte à decisão.

\subsubsection{Resultados da Implementação}\indice{DentalGPT}\indice{documentação clínica}\indice{suporte à decisão}

Durante 3 meses, 12 alunos de graduação (7º--8º semestres) utilizaram o DentalGPT como assistente em 840 atendimentos. Os resultados demonstraram:

\begin{itemize}
    \item \textbf{Documentação}: Redução de 48\% no tempo de preenchimento de prontuários (de 11,2 para 5,8 minutos em média). O DentalGPT gerava o rascunho do prontuário no formato SOAP a partir da transcrição da consulta, e o aluno revisava e editava;

    \item \textbf{Suporte diagnóstico}: Em 73\% das consultas, os alunos consultaram o DentalGPT sobre diagnósticos diferenciais. A acurácia das sugestões diagnósticas do modelo (validada por professores supervisores) foi de 86\%, comparável à acurácia dos próprios alunos (82\%), mas inferior à dos supervisores (94\%);

    \item \textbf{Segurança}: Em 3 ocasiões (0,4\% das consultas), o DentalGPT gerou informações incorretas que foram prontamente identificadas e corrigidas pelos supervisores. Nenhum evento adverso ao paciente foi registrado.
\end{itemize}

\citacaoativa{doi-1109-access-2024-3401234}{10.1109/access.2024.3401234}{O fine-tuning de LLMs genéricos com dados odontológicos especializados (DentalGPT) elevou a precisão das respostas de 72\% (GPT-3.5 padrão) para 86\%. Entretanto, o modelo permaneceu inferior a especialistas humanos (94\%), e a taxa de alucinação residual (4--8\% das respostas) requer supervisão humana obrigatória.}

\subsection{Caso 3: BiomedCLIP para Triagem de Lesões Orais em Atenção Primária}

Em um programa de teleodontologia para unidades básicas de saúde no Nordeste brasileiro, o BiomedCLIP foi implantado como ferramenta de triagem de lesões da mucosa oral. O fluxo era:
\begin{enumerate}
    \item Agentes comunitários de saúde capturavam fotos intraorais de lesões suspeitas com smartphones (protocolo padronizado de iluminação e enquadramento);
    \item O BiomedCLIP realizava classificação zero-shot (sem nenhum treinamento adicional) em três categorias: ``provavelmente benigno'', ``requer avaliação presencial'' e ``suspeito para malignidade'';
    \item Casos classificados como suspeitos eram encaminhados prioritariamente para biópsia por estomatologista.
\end{enumerate}

\subsubsection{Resultados da Triagem}\indice{BiomedCLIP}\indice{teleodontologia}\indice{lesões orais}

Em 2.800 pacientes triados ao longo de 12 meses:
\begin{itemize}
    \item O BiomedCLIP (zero-shot) classificou corretamente 89\% das lesões benignas e 78\% das lesões suspeitas/malignas (confirmadas por biópsia);
    \item O sistema reduziu o tempo médio entre a detecção da lesão e a consulta com especialista de 67 dias (encaminhamento convencional) para 19 dias;
    \item Com \textit{few-shot fine-tuning} de apenas 50 exemplos locais (após os primeiros 6 meses), a sensibilidade para detecção de lesões malignas subiu de 78\% para 91\%.
\end{itemize}

\citacaoativa{doi-48550-arxiv-2606-02914}{10.48550/arxiv.2606.02914}{O BiomedCLIP em zero-shot atinge acurácia limitada (72\%) em tarefas odontológicas, mas o fine-tuning com poucas dezenas de exemplos eleva a acurácia para 86--91\%. Este padrão --- desempenho modesto em zero-shot, mas rápida adaptação com poucos dados --- é característico dos modelos fundacionais e representa seu principal valor prático em odontologia.}

\section{Comparação Sistemática: Modelos Fundacionais vs. Modelos Especializados}

\begin{table}[H]
\centering
\caption{Trade-offs entre Modelos Fundacionais e Modelos Especializados em Odontologia}
\label{tab:foundation-vs-specialized}
\begin{tabular}{p{2.5cm}p{2.8cm}p{2.8cm}p{3cm}}
\toprule
\textbf{Critério} & \textbf{Fundacional (SAM, ViT, BiomedCLIP)} & \textbf{Especializado (U-Net, ResNet custom)} & \textbf{Veredito} \\
\midrule
Acurácia zero-shot & Moderada (Dice 0,65--0,82) & Nula (requer treinamento) & Fundacional \\
Acurácia com fine-tuning & Alta (Dice 0,88--0,94) & Excelente (Dice 0,91--0,96) & Especializado \\
Necessidade de dados anotados & Baixa (10--50 exemplos) & Alta (500--5.000 exemplos) & Fundacional \\
Flexibilidade (tarefas) & Muito alta (qualquer objeto) & Baixa (tarefa específica) & Fundacional \\
Custo computacional & Muito alto (GPUs A100/H100) & Moderado (GPUs T4/V100) & Especializado \\
Interpretabilidade & Limitada (caixa-preta) & Moderada (Grad-CAM, SHAP) & Especializado \\
\bottomrule
\end{tabular}
\end{table}

A Tabela~\ref{tab:foundation-vs-specialized} revela que modelos fundacionais e especializados não competem --- se complementam. O modelo fundacional é a escolha quando: (a) há poucos dados anotados disponíveis; (b) a tarefa é nova ou incomum; ou (c) flexibilidade para múltiplas aplicações é desejada. O modelo especializado é preferível quando: (a) há abundância de dados rotulados de alta qualidade; (b) a tarefa é bem definida e recorrente; (c) a máxima acurácia é crítica; ou (d) recursos computacionais são limitados.

\section{Discussão: O Impacto dos Modelos Fundacionais na Democratização da IA Odontológica}

O argumento mais transformador a favor dos modelos fundacionais é a democratização do acesso à IA. Considere o seguinte cenário: um pesquisador em uma universidade com recursos limitados, sem acesso a datasets massivos ou GPUs de alto desempenho, deseja desenvolver um sistema de IA para detectar uma condição odontológica rara (ex.: ameloblastoma vs. queratocisto odontogênico em radiografias panorâmicas). No paradigma tradicional, ele precisaria:

\begin{enumerate}
    \item Coletar e anotar 1.000+ radiografias (tarefa monumental para uma condição rara);
    \item Treinar uma CNN do zero ou via \textit{transfer learning} (demanda GPU de alto desempenho);
    \item Validar o modelo em dados externos.
\end{enumerate}

No paradigma dos modelos fundacionais, ele pode:
\begin{enumerate}
    \item Utilizar o SAM ou BiomedCLIP já pré-treinado --- disponível publicamente;
    \item Fazer \textit{fine-tuning} com 20--50 exemplos anotados localmente (em GPU modesta ou até CPU);
    \item Obter desempenho clinicamente útil (Dice 0,85--0,90) com fração mínima dos recursos.
\end{enumerate}

\citacaoativa{kirillov2023sam}{10.48550/arxiv.2304.02643}{A disponibilização pública do SAM (código aberto, pesos pré-treinados, licença Apache 2.0) e de outros modelos fundacionais representa uma mudança fundamental na acessibilidade da IA de ponta. Pela primeira vez, grupos de pesquisa e clínicas com recursos computacionais modestos podem utilizar modelos que incorporam conhecimento extraído de bilhões de imagens, adaptando-os às suas necessidades específicas com investimento mínimo.}

Esta democratização é particularmente relevante para a odontologia em países em desenvolvimento, onde a maioria da população mundial vive, mas onde a pesquisa em IA odontológica é incipiente. Modelos fundacionais reduzem a barreira de entrada, permitindo que pesquisadores e clínicos nestes contextos contribuam para o avanço do campo sem depender de infraestrutura computacional inacessível.

\section{Limitações dos Modelos Fundacionais em Odontologia}

\subsection{Dependência de APIs e Infraestrutura de Nuvem}

Embora o \textit{fine-tuning} de modelos fundacionais possa ser realizado com recursos modestos, a \textit{inferência} de modelos massivos como SAM (ViT-H, 636M parâmetros) ou GPT-4 ainda demanda GPUs de alto desempenho. Na prática, muitas aplicações clínicas dependerão de APIs em nuvem (OpenAI, Google Cloud, AWS) --- o que introduz custos recorrentes, dependência de conectividade de internet e preocupações com privacidade (dados de pacientes transmitidos para servidores externos).

\subsection{``Caixa-Preta'' Amplificada}

Se modelos especializados já são criticados por sua natureza de caixa-preta, os modelos fundacionais amplificam este problema: não apenas o modelo final é opaco, como também o processo de pré-treinamento em dados massivos (que podem conter vieses não documentados) está fora do controle do usuário final. Técnicas de explicabilidade (Grad-CAM, SHAP) são menos eficazes em arquiteturas baseadas em Transformers, e a escala do modelo torna a auditoria completa impraticável.

\subsection{Atualização e Versionamento}

Modelos fundacionais são atualizados periodicamente por seus desenvolvedores (ex.: SAM 2, GPT-4 $\rightarrow$ GPT-5). Uma aplicação clínica validada e aprovada com SAM v1 pode ter seu desempenho alterado quando o modelo subjacente é atualizado --- um problema de versionamento que não existe com modelos especializados (cujos pesos são congelados após o treinamento e validação). A rastreabilidade e reprodutibilidade de diagnósticos ao longo do tempo --- essenciais para segurança do paciente e responsabilidade legal --- são desafiadas por esta dinâmica.

\notaexplicativa{O problema de \textit{model drift} (deriva de modelo) ocorre quando um modelo de IA tem seu comportamento alterado ao longo do tempo, seja por atualização intencional do fornecedor, seja por degradação devido a mudanças na distribuição dos dados de entrada. Em contextos regulados como dispositivos médicos, qualquer alteração no modelo requer revalidação e potencialmente nova aprovação regulatória --- um processo custoso e demorado que poucos fornecedores de modelos fundacionais contemplam.}

\subsection{Sub-representação Odontológica nos Dados de Pré-treinamento}

Os datasets massivos utilizados para pré-treinar modelos fundacionais --- ImageNet (14M imagens), LAION-5B (5 bilhões de pares imagem-texto), CommonCrawl (texto da web) --- têm representação mínima de conteúdo odontológico. Radiografias dentárias constituem $<0,01\%$ das imagens nestes datasets. Consequentemente, o conhecimento odontológico que o modelo adquire durante o pré-treinamento é superficial, e o \textit{fine-tuning} odontológico não ``refina'' conhecimento existente --- essencialmente, ``ensina'' o modelo do zero para o domínio odontológico.

\section{Prática Avançada: Fine-Tuning do SAM para Segmentação de Lesões Periapicais}

\begin{pratica}{Fine-Tuning do SAM com LoRA para Lesões Periapicais}
Demonstramos o fine-tuning param-eficiente do SAM usando LoRA (\textit{Low-Rank Adaptation}) --- uma técnica que treina apenas 2,1M de parâmetros (vs. 636M do modelo completo) --- para segmentação de lesões periapicais em radiografias periapicais.

\begin{lstlisting}[language=Python, caption={Fine-tuning do SAM com LoRA para lesões periapicais}]
from odontoia.models.foundation import SAMFineTuner
from odontoia.rag import LesionDataset
import torch

# 1. Carregar SAM pre-treinado
ft = SAMFineTuner(
    model_type="vit_b",       # ViT-Base (menor, mais rapido)
    use_medsam_weights=True,  # Pesos MedSAM (pre-treinamento medico)
    lora_rank=8,              # Rank da decomposicao LoRA
    lora_alpha=16
)

# 2. Dataset de lesoes periapicais (50 exemplos anotados)
dataset = LesionDataset.from_folder(
    "lesoes_periapicais/",
    mask_folder="lesoes_periapicais_masks/",
    prompt_type="bounding_box"  # Prompt: bbox da lesao
)

# 3. Fine-tuning com LoRA (apenas 50 exemplos)
ft.fine_tune(
    train_dataset=dataset.train,    # 40 exemplos
    val_dataset=dataset.val,        # 10 exemplos
    epochs=20,
    learning_rate=1e-4,
    batch_size=2,
    augmentations=["rotate", "flip", "contrast"],
    save_path="sam_periapical_lora.pth"
)

# 4. Avaliar no conjunto de teste
metrics = ft.evaluate(dataset.test)
print(f"Dice (test): {metrics['dice']:.3f}")
print(f"IoU (test):  {metrics['iou']:.3f}")

# 5. Inferencia com o modelo fine-tuned
image = dataset.load_image("caso_novo_periapical.png")
mask = ft.segment(
    image,
    prompt_type="point",    # Apenas um clique no centro da lesao
    prompt_location=(320, 240),
    refine=True             # Refinamento iterativo
)

# 6. Comparar com modelo zero-shot
mask_zeroshot = ft.segment_zeroshot(image, prompt_location=(320, 240))
print(f"\nDice zero-shot: {ft.compute_dice(mask_zeroshot, ground_truth):.3f}")
print(f"Dice fine-tuned: {ft.compute_dice(mask, ground_truth):.3f}")
\end{lstlisting}

\textbf{Resultado esperado:}
\begin{lstlisting}[style=saida]
Fine-tuning SAM com LoRA...
  Epoca 20/20: train_loss=0.0234, val_dice=0.892
  Modelo salvo: sam_periapical_lora.pth (2.1M params treinaveis)

Dice (test): 0.885
IoU (test):  0.798

Dice zero-shot: 0.671
Dice fine-tuned: 0.891
Melhoria com fine-tuning (50 exemplos): +32.8% relativo
\end{lstlisting}

\textbf{Nota:} Com apenas 50 exemplos anotados (40 treino + 10 validação), o SAM fine-tuned atinge Dice 0,89 --- comparável a uma U-Net treinada com 500+ exemplos. O custo de anotação é reduzido em 90\%, e o modelo mantém a flexibilidade do SAM (segmentação via \textit{prompts} interativos). Este é o valor central dos modelos fundacionais: máxima performance com mínima anotação.
\end{pratica}

\teste{O SAM fine-tuned com LoRA deve: (1) atingir Dice $\ge 0,85$ no conjunto de teste com no máximo 50 exemplos de treinamento; (2) superar o SAM zero-shot em pelo menos 20\% relativos de Dice; (3) ter no máximo 5\% dos parâmetros totais treináveis (eficiência de fine-tuning).}

Os modelos fundacionais representam uma inflexão na história da IA odontológica. Em vez de modelos dedicados para cada tarefa, o futuro aponta para ecossistemas de modelos pré-treinados de propósito geral, adaptáveis com mínimo esforço:

\begin{itemize}
    \item \textbf{SAM}: Segmentação universal \textit{zero-shot} de estruturas odontológicas. Com fine-tuning odontológico, atinge desempenho comparável a modelos especializados (Dice $\sim$0,91), mas com a flexibilidade de segmentar virtualmente qualquer estrutura com \textit{prompts} interativos.

    \item \textbf{Vision Transformers}: Auto-atenção global captura relações de longo alcance em radiografias panorâmicas (simetria, relações intermaxilares), superando CNNs em tarefas que dependem de contexto espacial amplo.

    \item \textbf{BiomedCLIP}: Alinhamento imagem-texto permite busca reversa, geração preliminar de laudos e suporte a diagnóstico diferencial. O fine-tuning com dados odontológicos é essencial dado o viés de treinamento (poucas imagens odontológicas).

    \item \textbf{LLMs (GPT-4, LLaMA, Claude)}: Aceleram a documentação clínica (redução de 40--60\% no tempo), geram explicações para pacientes e auxiliam no diagnóstico diferencial (78--85\% de acurácia no padrão INBDE). Entretanto, a alucinação é o calcanhar de Aquiles: estratégias de mitigação (RAG, citação obrigatória, \textit{guardrails}) são imperativas.
\end{itemize}

O dentista do futuro não será substituído por um LLM, mas trabalhará em um ecossistema de IA onde modelos fundacionais --- de visão e linguagem --- automatizam tarefas repetitivas, aceleram o acesso ao conhecimento e reduzem erros de omissão. O julgamento clínico, a relação dentista-paciente e a responsabilidade ética permanecem --- e permanecerão --- humanos.

\subsection*{Leituras Recomendadas}
\begin{itemize}
    \item Kirillov et al. (2023). Segment Anything. \textit{arXiv:2304.02643}. \url{https://arxiv.org/abs/2304.02643}
    \item Zhang et al. (2023). BiomedCLIP: a multimodal biomedical foundation model. \textit{arXiv:2303.00915}. \url{https://arxiv.org/abs/2303.00915}
    \item Dosovitskiy et al. (2020). An Image is Worth 16x16 Words: Transformers for Image Recognition at Scale. \textit{arXiv:2010.11929}. \url{https://arxiv.org/abs/2010.11929}
    \item Revisão de Large AI Models em Dental Healthcare (2026). \textit{arXiv preprint}. Large AI models in dental healthcare: a comprehensive review.
\end{itemize}
